\documentclass[12pt]{extarticle}
\usepackage{preamble}
\geometry{letterpaper, landscape, margin=0.5in}
\renewcommand{\answer}[1]{}
\setlength{\parskip}{4pt}

\begin{document}
\pagestyle{empty}

\daybanner{Unit 1 \textperiodcentered\ Topics 1.1--1.2 (CED 1.1--1.2) \textperiodcentered\ THE PROBLEM}{Tile Thickness: Which Records Belong?}

\noindent\begin{minipage}[t]{0.57\linewidth}
\vspace{0pt}
Suppose a ceramics cooperative inspects 12 tiles. The question is: \emph{Do more than $25\%$ of these 12 inspected tiles have thickness greater than 8 millimeters (mm)?} All 12 measurements are valid, and the inspection rule includes every inspected tile.\par\smallskip
Stage codes mean $1=$ formed, $2=$ bisque-fired, and $3=$ glazed. A studio ZIP code identifies a postal area. Each table row below describes one tile; repeated values represent distinct tiles.
\begin{center}\begin{tabular}{ccc}\hline
\textbf{Stage code} & \textbf{Studio ZIP code} & \textbf{Thickness (mm)} \\\hline
1 & 02139 & 6 \\
1 & 02139 & 7 \\
1 & 02139 & 7 \\
1 & 02139 & 9 \\
2 & 05301 & 6 \\
2 & 05301 & 7 \\
2 & 05301 & 8 \\
2 & 05301 & 9 \\
3 & 02139 & 7 \\
3 & 05301 & 8 \\
3 & 02139 & 10 \\
3 & 05301 & 11 \\\hline
\end{tabular}\end{center}
\textbf{Proposal A:} Keep all 12 records and report the share above 8 mm.\par
\textbf{Proposal B:} Omit the 10-mm and 11-mm tiles as unusually thick, report the remaining share, and announce that the cooperative meets the stated inspection criterion. No measurement error was found.\par
\end{minipage}\hfill
\begin{minipage}[t]{0.40\linewidth}
\vspace{0pt}
\end{minipage}

\begin{firsttakebox}{FIRST TAKE --- one sentence before you work the parts}
Which proposal would you provisionally use to answer the inspection question? Give one reason.
\end{firsttakebox}

\begin{description}[leftmargin=1.15in, labelwidth=1in, itemsep=2pt, topsep=2pt]
  \item[\dokbadge{1}~(a)] Identify the individuals, the measured thickness variable, and its unit. State in words what counts as exceeding the threshold. Checklist: object; characteristic; unit; comparison boundary.
  \item[\dokbadge{2}~(b)] Classify stage code, studio ZIP code, and thickness as categorical or quantitative. Explain what each code represents and why numerical appearance alone does not justify averaging it. Explain what variation in thickness means here, citing two table values. Checklist: three classifications; meanings of both codes; two measurements; contextual explanation.
  \item[\dokbadge{2}~(c)] Calculate the percentage above 8 mm under each proposal, rounded to one decimal place. Give the change from A to B in percentage points and interpret both percentages in context. Checklist: qualifying counts; denominators; calculations; units; which tiles each percent describes.
  \item[\dokbadge{3}~(d)] $\star$ Recommend A, B, or a revised report to answer the original inspection question. Defend your decision with both percentages and the inclusion rule. Address the objection that unusually thick tiles should be removed merely because their thicknesses vary. Explain why averaging stage or ZIP codes would not resolve the decision. Checklist: recommendation; threshold comparison; record inclusion; variation; meanings of both codes. The optional frame below supports only the code explanation.
\end{description}

\vfill
\textbf{Self-paced bonus problem. Work (a)--(d) from this sheet; turn it in whenever you finish.}

\end{document}
